\documentclass[11pt]{article}
\usepackage[T1]{fontenc}
\usepackage{lmodern,amsmath,amssymb,amsthm}
\usepackage[margin=1in]{geometry}
\usepackage[hidelinks]{hyperref}
\newtheorem{theorem}{Theorem}
\newtheorem{proposition}[theorem]{Proposition}
\newtheorem{lemma}[theorem]{Lemma}
\DeclareMathOperator{\Cl}{Cl}
\newcommand{\leg}[2]{\left(\frac{#1}{#2}\right)}
\title{A conductor correction for class-number tests of twin-prime neighbors}
\author{\small \href{mailto:admin@cosmosgeek.org}{admin@cosmosgeek.org}\qquad \href{mailto:b321014@iiit-bh.ac.in}{b321014@iiit-bh.ac.in}}
\date{Public research draft, November 2026}
\begin{document}
\maketitle
\begin{abstract}
For an even midpoint $m\ge4$, two congruences of imaginary quadratic
order class numbers, together with two auxiliary Jacobi symbols, are
equivalent to the simultaneous primality of $m-1$ and $m+1$. One
congruence removes unramified conductor primes; a square-divisor
M\"obius difference removes conductors supported on ramified primes.
For each fixed auxiliary prime $\ell\equiv1\bmod4$, a Pell family gives
infinitely many even midpoints with both neighbors composite and
Hurwitz class number $H(\ell(m^2-1))\equiv4\bmod8$.
A second calculation gives the asymptotic count of fundamental
discriminants passing the test when the auxiliary prime varies.
For primes $\ell\equiv37\pmod {140}$ we also prove that
$6\ell-1,6\ell+1$ are both prime exactly when the number of primitive
ordered, signed three-square representations of $36\ell^2-1$ is
$48\bmod96$. We give an infinite family of composite neighbors for
which a square-index mock-theta coefficient takes the twin residue.
A second family shows that a nonzero class-number trace at the
correct square index can come entirely from other trace coordinates.
\end{abstract}

\section{The two order class numbers}
For $-n<0$ a discriminant, let $h(-n)$ denote the proper ideal class
number of the order of discriminant $-n$. For the odd integers
$n\equiv3\bmod4$ occurring below put
\[
 J(n)=\sum_{s^2\mid n}\mu(s)\,h(-n/s^2).
\]
All the arguments $-n/s^2$ here are again discriminants.
For even $m\ge4$ put $a=m-1$ and $b=m+1$. Call a prime
$\ell\equiv1\bmod4$ admissible for $m$ if its Jacobi symbols satisfy
\[
             \leg{\ell}{a}=\leg{\ell}{b}=-1 .
\]
In particular $\ell\nmid ab$.

\begin{theorem}\label{th:main}
The integers $m-1$ and $m+1$ are both prime if and only if there
is an admissible prime $\ell$ for which, with $n=\ell(m^2-1)$,
\[
                  h(-n)\equiv J(n)\equiv4\pmod8.
\]
If both neighbors are prime, every admissible $\ell$ works, and there
are infinitely many such $\ell$.
\end{theorem}

We use two classical facts. For a negative fundamental discriminant
with $t$ distinct prime-discriminant factors, the $2$-rank of its
class group is $t-1$, and its $4$-rank is $t-1$ minus the rank of
the R\'edei matrix over $\mathbb F_2$ \cite[\S\S2--3]{stev}.
For a positive odd squarefree $d\equiv3\bmod4$ with $d>3$, the ring
class number formula \cite[Theorem~7.24]{cox} is
\begin{equation}\label{eq:order}
 h(-df^2)=h(-d) A_d(f),\qquad
 A_d(f)=f\prod_{t\mid f,\ t\ {\rm prime}}
 \left(1-\frac{\chi_{-d}(t)}{t}\right),
\end{equation}
where $\chi_{-d}$ is the primitive quadratic character. Both the
maximal order and the orders in \eqref{eq:order} have just two units.

\begin{lemma}\label{lem:redei}
Let $p,q,\ell$ be distinct odd primes, exactly one of $p,q$ congruent to
$3\bmod4$, and $\ell\equiv1\bmod4$. Let $x,y,z$ be the additive
Legendre symbols of $(q/p),(\ell/p),(\ell/q)$, respectively, with a
nonresidue encoded by $1$. Then
\[
               h(-\ell pq)\equiv4\bmod8
               \quad\Longleftrightarrow\quad xy+xz+yz=1
               \quad\text{in }\mathbb F_2.
\]
\end{lemma}
\begin{proof}
The discriminant $-\ell pq$ is fundamental. The off-diagonal entries
of its R\'edei matrix are symmetric by quadratic reciprocity: the
single negative prime-discriminant factor is paired only with primes
$1\bmod4$. With the diagonal determined by zero column sums, the matrix is
\[
 \begin{pmatrix}x+y&x&y\\x&x+z&z\\y&z&y+z\end{pmatrix}.
\]
Its rank is two precisely when the displayed quadratic polynomial is one:
all cofactors are equal, since the row and column sums vanish.
The $2$-rank is two. The $2$-primary part therefore has order four
exactly when the $4$-rank is zero.
\end{proof}

\begin{proof}[Proof of Theorem~\ref{th:main}]
Write $n=df^2$ with $d$ squarefree, and fix an admissible $\ell$.
Since $ab\equiv3\bmod4$, $-d$ is fundamental. The coprime odd integers
$a$ and $b$ each have a prime of odd valuation, by the two negative
Jacobi symbols. Together with $\ell$ this gives at least three prime
factors of $d$, so $4\mid h(-d)$ and $d>3$.

The function $A_d$ in \eqref{eq:order} is multiplicative and
$A_d(t^e)=t^{e-1}(t-\chi_{-d}(t))$ for $e\ge1$. If an unramified
prime divides $f$, this factor is even and $8\mid h(-n)$.
If $f>1$ is supported entirely on ramified primes, then $A_d(c)=c$
for every divisor $c$ of $f$, so
\[
 J(n)=h(-d)\sum_{s\mid f}\mu(s)A_d(f/s)
     =h(-d)\varphi(f)\equiv0\pmod8.
\]
Here $f$ is odd. Thus joint success forces $f=1$. A composite neighbor
would now supply a fourth distinct prime factor of $d$, forcing
$8\mid h(-d)$ by genus theory.

Conversely, for prime neighbors $p=a,q=b$, use the Chinese remainder
theorem and Dirichlet's theorem to choose $\ell\equiv1\bmod4$ a
nonresidue modulo both. In Lemma~\ref{lem:redei}, $y=z=1$ and
$xy+xz+yz=1$ regardless of $x$. We have $f=1$, so $J(n)=h(-n)$.
\end{proof}

\section{How many discriminants can come from fixed neighbors?}
On the squarefree subfamily there is also an exact count without
imposing the two negative-symbol conditions. For a fixed genuine twin
center, $x$ in Lemma~\ref{lem:redei} is determined by $(2/(m-1))$.
Among primes $\ell\equiv1\bmod4$, the four pairs $(y,z)$ have equal
limiting frequency by the prime number theorem in fixed progressions.
If $x=0$ only $(1,1)$ works; if $x=1$ all but $(0,0)$ work.

\begin{proposition}
Let $N(X)$ count distinct fundamental discriminants
$-\ell(m^2-1)$ of absolute value at most $X$ with $m\ge4$ even,
$\ell>2m$ prime congruent to $1\bmod4$, and class number $4\bmod8$.
Then
\[
 N(X)\sim\frac X{\log X}
 \sum_{\substack{m\ge4\text{ even}\\m-1,m+1\text{ prime}}}
 \frac{\delta_m}{m^2-1},
 \qquad
 \delta_m=
 \begin{cases}
 1/8,&m\equiv0,2\pmod8,\\
 3/8,&m\equiv4,6\pmod8 .
 \end{cases}
\]
\end{proposition}
\begin{proof}
Genus theory forces both neighbors to be prime. The largest prime
factor is $\ell$, so the representation is unique. The calculated
densities, including the density $1/2$ for $\ell\equiv1\bmod4$, give
the formula for any finite sum over $m$. For $B<m\le X^{1/4}$,
$\pi(X/(m^2-1))\ll X/(m^2\log X)$ and the sum is
$O(X/(B\log X))$. For larger $m$ the trivial bound gives
$O(X^{3/4})$. Let $X\to\infty$ and then $B\to\infty$.
The single center $m=4$ contributes $1/40$ to the sum.
\end{proof}
In particular the main term is compatible with finitely many genuine
centers. A result forcing arbitrarily large successful midpoints is
still needed.

\section{Composite neighbors in a Hurwitz congruence}
Let $H(n)$ be the Hurwitz class number, normalized by weights
$1/|\mathrm{Aut}_{\mathrm{PSL}_2}(Q)|$. For $d$ as in
\eqref{eq:order}, partitioning forms by their common content gives
\begin{equation}\label{eq:hurwitz}
 H(df^2)=\sum_{c\mid f}h(-dc^2)
        =h(-d)\prod_{t^e\parallel f}
                 \left(1+\sum_{j=1}^e A_d(t^j)\right).
\end{equation}
Here and below $t^e\parallel f$ in a product denotes exact prime
powers. This is the usual conductor formula for Hurwitz coefficients
\cite{brr}. For odd $t\nmid d$ the corresponding factor is odd for
every $e$. For odd $t\mid d$ its parity is $e+1\bmod2$.

\begin{theorem}\label{th:pell}
For every prime $\ell\equiv1\bmod4$, infinitely many even $m$
have two composite neighbors and $H(\ell(m^2-1))\equiv4\bmod8$.
For $\ell=13$ this holds even with $\ell$ admissible for $m$.
\end{theorem}
\begin{proof}
Choose primes $u,v>3$ distinct from $\ell$, with
\[
 u\equiv5\bmod8,\quad v\equiv3\bmod4,\quad
 \leg{\ell}{u}=\leg{\ell}{v}=\leg{u}{v}=-1 .
\]
To do this, choose $u$ in a nonresidue class modulo $\ell$ and the
specified class modulo eight; then impose independent nonresidue
classes modulo $\ell,u$ in choosing $v$. Reciprocity, CRT and
Dirichlet's theorem give the required primes. Set $D=uv$.
Lemma~\ref{lem:redei} yields $h(-\ell D)\equiv4\bmod8$.

Let $x+y\sqrt D>1$ be the least norm-one positive Pell unit in
$\mathbb Z[\sqrt D]$; $x,y>0$. We claim $x$ is even. If $x$ were odd,
$y$ would be even and the coprime consecutive integers
$(x+1)/2,(x-1)/2$ would have product $D(y/2)^2$.
If the first is a square and the second is $D$ times a square, the
Pell unit is the square of a smaller positive norm-one unit.
The reversed allocation would solve $s^2-Dt^2=-1$, impossible
modulo $v\equiv3\bmod4$. The remaining allocations would give
$u s^2-vt^2=1$ or $v s^2-ut^2=1$. Reduction modulo $v$ or $u$,
respectively, contradicts $(u/v)=(v/u)=-1$. This proves the claim.

Every odd power $x'+y'\sqrt D$ of this unit has $x'$ even, $y'$ odd
(by reduction modulo two). Each of the coprime odd neighbors $x'\pm1$
must carry one of the two squarefree factors $u,v$ in
$(x'-1)(x'+1)=D(y')^2$. Indeed, if both factors belonged to one side,
the other would be a square congruent to $2$ or $-2$ modulo $u$.
Both are nonresidues since $u\equiv5\bmod8$. Thus for large $x'$ both
neighbors are composite.

We arrange that $v_t(y')$ is even at $t=u,v,\ell$. For any Pell unit
$X+Y\sqrt D$ and prime $t\ge5$ dividing $DY^2$, its $n$th power
$X_n+Y_n\sqrt D$ satisfies
\[
 v_t(Y_n)=v_t(Y)+v_t(n).
\]
Indeed, in $Y_n/(nY)$ the first binomial term is $X^{n-1}$,
a $t$-adic unit. The terms indexed by $j\ge1$ have factors
$(DY^2)^j\binom{n-1}{2j}/(2j+1)$ and are divisible by $t$, since
$j>v_t(2j+1)$. First take an odd power with an additional factor
$u$ or $v$ precisely for each of these two primes where the valuation
is odd. In the resulting unit, if $v_\ell(Y)$ is odd, take the
$\ell$th power. This makes all three valuations even.

We may now raise the resulting unit to arbitrarily large odd powers
$n$ with $(n,uv\ell)=1$ and $n\equiv1$ modulo its multiplicative
order in $(\mathbb F_\ell[\sqrt D])^\times$ (the quadratic algebra
here means $\mathbb F_\ell[T]/(T^2-D)$). This also preserves
$\ell\nmid Y_n$ when $v_\ell(Y)=0$. All three valuations stay even.
Formula \eqref{eq:hurwitz} for $d=\ell D$, $f=Y_n$ gives the required
congruence at $m=X_n$.

For the stronger example start directly with $6+\sqrt{35}$ and
$\ell=13$. Take arbitrarily large powers congruent to one modulo
the order of this unit in
$((\mathbb Z/(2\cdot5\cdot7\cdot13))[T]/(T^2-35))^\times$.
The new $m$ is even, $Y_n$ is coprime to $2\cdot5\cdot7\cdot13$,
and $m-1,m+1$ have squarefree parts $5,7$ respectively (the chosen
congruences fix their allocation). Both are composite for $m>8$.
Since $(13/5)=(13/7)=-1$, both Jacobi gates hold.
\end{proof}

\section{A primitive three-square criterion}
For $D>0$, let $r_3(D)$ count all ordered triples of integers
$(x,y,z)$ satisfying $x^2+y^2+z^2=D$, with every choice of signs
counted. Let $r_3^*(D)$ count those with $\gcd(x,y,z)=1$. In this
section $m>2$ and $m\equiv6\bmod8$. Write
\[
 D=m^2-1=(m-1)(m+1)=d f^2,\qquad d\ \text{squarefree}.
\]
The neighbors are coprime and have residues $5,7\bmod8$. Neither is
a square, so each contributes a different prime to $d$. In particular
$d>3$, $d\equiv3\bmod8$, and $-d$ is fundamental. Every order in the
field has exactly two units. This removes the exceptional automorphism
weights at discriminants $-3,-4$.

\begin{lemma}\label{lem:primitive}
For these midpoints,
\[
 \frac{r_3^*(D)}{24}=h(-D),\qquad
 h(-D)\equiv2\bmod4
 \ \Longleftrightarrow\
 m-1=p^{2a+1},\quad m+1=q^{2b+1},
\]
for distinct primes $p,q$ and integers $a,b\ge0$.
\end{lemma}
\begin{proof}
Gauss's three-square formula gives $r_3(E)=24H(E)$ for every
$E\equiv3\bmod8$ \cite[Theorem 1.4]{mortenson}. The primitive
formula is itself classical \cite[Theorem 1.3]{mortenson}; here the
square-divisor inversion also identifies the conductor correction.
Each $g^2\mid D$
has $g\mid f$ and $D/g^2\equiv3\bmod8$. Removing the common gcd of
the coordinates and using \eqref{eq:hurwitz}, we obtain
\[
 r_3^*(D)/24=\sum_{g\mid f}\mu(g)H(D/g^2)=h(-df^2).
\]
Genus theory gives $2^{\omega(d)-1}\mid h(-d)$, where $\omega$ counts
distinct prime factors. In \eqref{eq:order}, a prime of $f$ outside
$d$ contributes the even factor $t^{e-1}(t-\chi_{-d}(t))$.
Consequently $h(-D)\equiv2\bmod4$ forces $\omega(d)=2$ and
every prime of $f$ to divide $d$. Coprimality of the neighbors
puts one base on each side, with an odd exponent on each.

Conversely these prime powers have $p\equiv5$, $q\equiv7\bmod8$.
Since $q^{2b+1}\equiv2\bmod p$, the Legendre symbol
$(q/p)=(2/p)=-1$. The squarefree two-prime class-number
criterion gives $h(-pq)\equiv2\bmod4$
\cite[Theorem 2.3(2), including its attribution to Hasse]{chen}.
Here $H(pq)=h(-pq)$ and $3H(pq)\equiv2$ is equivalent to
$h(-pq)\equiv2\bmod4$. Every conductor factor is now a power of a
ramified odd prime, hence odd. Equation \eqref{eq:order} proves
the assertion for $D$.
\end{proof}

\begin{theorem}\label{th:three}
For every prime $\ell\equiv37\bmod140$,
\[
       r_3^*(36\ell^2-1)\equiv48\bmod96
       \quad\Longleftrightarrow\quad
       6\ell-1,\;6\ell+1\text{ are both prime}.
\]
\end{theorem}
\begin{proof}
The hypotheses give $\ell\equiv1\bmod4$ and $\ell\equiv2\bmod35$.
Take $m=6\ell$ in Lemma~\ref{lem:primitive}. If $6\ell-1=p^k$
with odd $k\ge3$, then $p+1$ divides $6\ell$ and its cofactor
$(p^k+1)/(p+1)$ is at least $p^2-p+1>3$. The even divisors of
$6\ell$ are $2,6,2\ell,6\ell$; $p\equiv5\bmod8$ leaves only
$p+1=6$. But $6\ell-1\equiv1\bmod5$. Similarly, if $6\ell+1=q^k$
with odd $k\ge3$, then $q-1$ divides $6\ell$ and its cofactor
is at least $q^2+q+1>3$. As $q\equiv7\bmod8$, this forces
$q-1=6$, whereas $6\ell+1\equiv6\bmod7$. Thus both exponents
equal one. The converse follows directly from the lemma.
\end{proof}

\section{A mock-theta coefficient at square indices}
Let $v(n)=[z^n]V_1(z)$ for the eighth-order mock theta function
\[
 V_1(z)=\sum_{j\ge0}\frac{z^{(j+1)^2}(-z;z^2)_j}{(z;z^2)_{j+1}},
 \qquad (a;b)_j=\prod_{i=0}^{j-1}(1-ab^i).
\]
Chen and Garvan prove $v(n)\equiv\chi(n)H(4n-1)\bmod4$ for $n>0$,
where $\chi(n)=-1$ on $0,1\bmod4$ and $+1$ otherwise
\cite[Lemma 3.4]{chen}. They also characterize $3H(D)\equiv2\bmod4$
for all $D\equiv3\bmod4$ in terms of two specified primes to exponents
$1\bmod4$ times a square \cite[Theorem 2.5]{chen}. These are
established inputs. Taking $n=(m/2)^2$ picks out the project
discriminant $4n-1=m^2-1$.

\begin{proposition}\label{prop:mock}
There exist infinitely many $m\equiv6\bmod8$ with
\[
   v\bigl((m/2)^2\bigr)\equiv2\bmod4,\qquad
   r_3^*(m^2-1)\equiv0\bmod96
\]
and with both $m-1$ and $m+1$ composite.
\end{proposition}
\begin{proof}
The matrix $\left(\begin{smallmatrix}6&7\\5&6\end{smallmatrix}\right)$
has determinant one and preserves $7v^2-5u^2=2$ on vectors $(u,v)$.
Its image of $(1,1)$ is $(13,11)$. Iterate a positive power congruent
to the identity modulo $70$, starting at $(13,11)$. This gives
unbounded positive odd solutions with $\gcd(uv,35)=1$. Put
$m=5u^2+1=7v^2-1$. Both neighbors are composite and
$D=35(uv)^2$. Since $(7/5)=-1$, $h(-35)\equiv2\bmod4$.
Every factor in the Hurwitz product \eqref{eq:hurwitz} is odd:
at $t\mid uv$ all the terms after its initial 1 are even. Hence
$H(D)\equiv2\bmod4$, and Chen--Garvan's congruence gives the asserted
value of the coefficient. The conductor $uv>1$ has an unramified
prime, so $4\mid h(-D)$ by \eqref{eq:order}. Apply
Lemma~\ref{lem:primitive}.
\end{proof}
For example $m=846$ gives $m-1=5\cdot13^2$, $m+1=7\cdot11^2$
and $(m/2)^2=178929$. At discriminant $-35$ the two reduced forms
are $[1,1,9]$ and $[3,1,3]$. The character $\chi_{35}$ equals $1$
at both $11$ and $13$. Thus the ordinary and primitive class
counts are respectively $2(1+10)(1+12)=286$ and $2\cdot10\cdot12=240$.
The proposition permits general even midpoints
in the indicated class. Its construction imposes no primality
condition on $m/6$.

A counterexample also occurs on the restricted prime-center sequence
itself. Take $\ell=6197\equiv37\bmod140$. Then
\[
 6\ell-1=37181,\qquad 6\ell+1=37183=19^2\cdot103 .
\]
Here $\ell$ and $37181$ are prime (trial division through $78$ and
$192$, respectively, suffices), as are $19$ and $103$. The squarefree
core is $d=37181\cdot103$, with $(103/37181)=(2/37181)=-1$ by the
displayed factorization. Thus $h(-d)\equiv2\bmod4$. For the sole
conductor prime $19\nmid d$, $1+A_d(19)$ is odd while $A_d(19)$ is
even. The square-index coefficient $v(345625281)$ consequently equals
$2\bmod4$ although the primitive count is $0\bmod96$. This gives
a single restricted-center example; Proposition~\ref{prop:mock}
gives an infinite family at general centers $m\equiv6\bmod8$.

The conductor correction can itself be expressed in these coefficients.
For every $m$ of Lemma~\ref{lem:primitive},
\[
   \sum_{g^2\mid D}\mu(g)\,
        v\!\left(\frac{D/g^2+1}{4}\right)
         \equiv r_3^*(D)/24\pmod4 .
\]
Indeed every $D/g^2$ has the same fundamental core $d$, hence its
Hurwitz class number is an even integer by \eqref{eq:hurwitz} and genus
theory. The sign $\chi$ therefore disappears modulo four. The
square-divisor inversion in Lemma~\ref{lem:primitive} gives the
formula, explaining exactly what is lost when only $g=1$ is used.

\section{What the square-index traces detect}
Let $N=9\ell^2$ with $\ell>7$ prime and $\ell\equiv1\bmod4$.
Take $H(0)=-1/12$, with $H(u)=0$ when $u<0$ or $u\equiv1,2\bmod4$.
Bringmann--Kane study sums of these coefficients on
fixed trace progressions \cite[Theorem 1.4]{bk}. Here an evaluable
second moment is
\[
 B(N)=\sum_{\substack{|t|\le2\sqrt N\\t\equiv1\ (3)}}t^2H(4N-t^2)
   =9\ell^2(9\ell^2+8\ell-1)-730-\ell^3
   \equiv1\bmod4 .
\]
For the equality use Bringmann--Kane--Pujahari
\cite[Theorem 1.1(2), nonzero residue class]{bkp}. The divisors of
$N$ strictly below $\sqrt N$ are $1,3,9,\ell$. Omitting those with
$3\mid d$ and $3\mid N/d$ leaves the cubic sum $1+9^3+\ell^3$.
The divisor-sum difference $\sigma(N)-\sigma(N/3)$ is
$9(1+\ell+\ell^2)$. In the sum of $\min(d,N/d)$ the central divisor
$3\ell$ cancels against the theorem's square-index correction; the
strictly smaller divisors have sum $13+\ell$. The remaining sum over
strictly smaller divisors of $N/9=\ell^2$ contributes $3N$.
Together these terms give the displayed polynomial. All class numbers
in this trace are integral: $4N-t^2\equiv2\bmod3$ excludes both
exceptional square classes $3u^2,4u^2$.

There is a related trace with a different primitivity condition.
For odd $s$ and $m=2s$, let $F_m(t)$ count positive odd ordered
triples $x,y,z$ with $t^2+x^2+y^2+z^2=m^2$ and
$\gcd(t,x,y,z)=1$. For odd $0<t<m$,
\begin{equation}\label{eq:fourmarked}
 F_m(t)=3\sum_{g\mid(t,s)}\mu(g)
            H\bigl((m^2-t^2)/g^2\bigr).
\end{equation}
In fact $m^2-t^2\equiv3\bmod8$ forces three odd coordinates.
Dividing Gauss's count by the eight sign choices and removing
the common four-coordinate gcd gives \eqref{eq:fourmarked}.

\begin{proposition}\label{prop:trace}
For $\ell$ as above and $m=6\ell$ one has
\[
   \sum_{\substack{0<t<m\\t\ {\rm odd},\ 3\nmid t}}F_m(t)
        =9\ell(\ell+1)\equiv2\bmod4 .
\]
For infinitely many primes $\ell\equiv37\bmod140$, this congruence
holds while both $F_m(1)$ and $r_3^*(m^2-1)/24$ are $0\bmod4$.
\end{proposition}
\begin{proof}
Write $\sigma(n)=\sum_{d\mid n}d$. Jacobi's formula
\cite{jacobi} gives $r_4(4s^2)=24\sigma(s^2)$ and
$r_4(s^2)=8\sigma(s^2)$. Subtracting the all-even quadruples
leaves all-odd quadruples by reduction modulo four. Dividing by
16 independent signs gives exactly $\sigma(s^2)$ positive ordered
odd quadruples. Their common gcd divides $s$. At $s=3\ell$ their
primitive count is therefore
\[
 \sum_{g\mid3\ell}\mu(g)\sigma((3\ell/g)^2)
  =(\sigma(9)-1)(\sigma(\ell^2)-1)=12\ell(\ell+1).
\]
Modulo three, a primitive quadruple of this norm has exactly three
coordinates prime to three. By permutation symmetry three quarters
have first coordinate prime to three, giving the trace.

Impose also $\ell\equiv2\bmod121$, compatible with the specified
reduced class modulo $140$. Dirichlet's theorem supplies infinitely
many primes. Then $11\parallel m-1$, and the squarefree kernel of
$m-1\equiv5\bmod8$ cannot equal $11\equiv3\bmod8$. The other
neighbor is a nonsquare coprime to the first. Hence $d$ has at
least three prime factors, and $4\mid h(-d)$. Equations
\eqref{eq:order} and \eqref{eq:hurwitz} give $4\mid h(-D),H(D)$.
At $t=1$ condition $\gcd(t,x,y,z)=1$ is automatic, so that
$F_m(1)=3H(D)$. This proves both vanishings. In particular the
other coordinates alone account for the residue $2$ of the trace.
\end{proof}
Already $\ell=37$ gives $m=222$ and $D=13\cdot17\cdot223$.
Its two marked residues vanish by genus theory while the two
complete traces above are respectively $1$ and $2$ modulo four.
For instance the displayed evaluations give
$B(12321)=155390353$ and $9\ell(\ell+1)=12654$.
At the same midpoint $(t,x,y,z)=(23,217,21,35)$ has norm $222^2$.
Its four-coordinate gcd is one and its last three coordinates
have gcd seven. This also exhibits the difference between the two
primitive counts on a single representation.

\begin{proposition}\label{prop:fixedtrace}
For every fixed trace modulus $M\ge1$ there are infinitely many
$m\equiv6\bmod8$ with $h(-(m^2-1))\equiv0\bmod4$ and another
odd trace $1<t<m$, $t\equiv1\bmod M$, for which
$m^2-t^2$ is squarefree and $H(m^2-t^2)\equiv2\bmod4$.
\end{proposition}
\begin{proof}
Let $M_0=[8,M]$ denote the least common multiple. Choose a prime
$r\equiv3\bmod8$ coprime to $M_0$. Choose a prime $p\equiv5\bmod8$
with $p\equiv1\bmod r$ and $\gcd(p(p+2),M_0)=1$.
At each odd prime dividing $M_0$ one can avoid $0,-2$; the Chinese
remainder theorem and Dirichlet's theorem give such $p$. For infinitely
many primes $q>p+2$ impose
\[
 q\equiv p+2\bmod{2M_0},\qquad q\equiv2\bmod p,\qquad
 q\equiv2-p+2r\bmod{r^2}.
\]
The moduli are coprime and the residues are reduced, the last being
$1\bmod r$. Set $m=(q+p)/2$, $t=(q-p)/2$. Then $t\equiv1\bmod M_0$,
$m\equiv6\bmod8$, and $r\parallel m-1$. The preceding
squarefree-kernel argument with $r$ in place of $11$ gives
$4\mid h(-(m^2-1))$. But $m^2-t^2=pq$ with $(q/p)=(2/p)=-1$,
so $H(pq)\equiv2\bmod4$ by \cite[Theorem 2.3]{chen}.
\end{proof}
The last proposition concerns termwise isolation at general midpoints.
It leaves open possible cancellations in the complete trace and the
distribution of individual primitive coefficients on the prime centers
of Theorem~\ref{th:three}.

\paragraph{Scope and provenance.}
The two-congruence test combines genus theory, the R\'edei matrix
and the ring class number formula. The Pell construction and the
fixed-center count describe the two distinct multiplicities a
nonvanishing argument must control. They give no unbounded sequence
of genuine twin centers. The original Twin Prime research chats and
Codex contributed the mathematical development, extraction, and
internal checks. Specialist review and the priority of these
specializations remain open. The same project's original class-number
chat supplied the primitive-midpoint criterion and the trace
counterexamples. The mock-theta specialization follows here from
Chen and Garvan's published coefficient theorem and the recovered Pell
construction. Codex checked the arithmetic and prepared the extension.
The recovered derivation is available in the
\href{https://github.com/cosmosgeek0/Twin-Prime/blob/main/history/hosted-proofs-supplement.md}{project source supplement, part 1}.

\begin{thebibliography}{9}
\bibitem{stev} P. Stevenhagen, \emph{R\'edei reciprocity, governing
fields, and negative Pell}, \S\S2--3, Theorem 3.1.
\href{https://arxiv.org/abs/1806.06250}{arXiv:1806.06250}.
\bibitem{cox} D. A. Cox, \emph{Primes of the Form $x^2+ny^2$},
Theorem 7.24.
\bibitem{brr} O. Beckwith, M. Raum and O. K. Richter,
\emph{Nonholomorphic Ramanujan-type congruences for Hurwitz class numbers}.
\href{https://doi.org/10.1073/pnas.2005984117}{Proc. Natl. Acad. Sci. USA 117 (2020)}.
\bibitem{chen} R. Chen and F. G. Garvan, \emph{A proof of the mod 4
unimodal sequence conjectures and related mock theta functions},
Theorems 2.3, 2.5 and Lemma 3.4.
\href{https://arxiv.org/abs/2010.14315}{arXiv:2010.14315}.
\bibitem{mortenson} E. T. Mortenson, \emph{A Kronecker-type identity
and the representations of a number as a sum of three squares},
Bull. Lond. Math. Soc. 49 (2017), 770--783.
\href{https://doi.org/10.1112/blms.12068}{doi:10.1112/blms.12068}.
\bibitem{jacobi} G. E. Andrews, S. B. Ekhad and D. Zeilberger,
\emph{A short proof of Jacobi's formula for the number of
representations of an integer as a sum of four squares},
Amer. Math. Monthly 100 (1993), 273--276.
\href{https://arxiv.org/abs/math/9206203}{arXiv:math/9206203}.
\bibitem{bk} K. Bringmann and B. Kane, \emph{Sums of class numbers
and mixed mock modular forms}, Theorem 1.4.
\href{https://doi.org/10.1017/S0305004118000348}{doi:10.1017/S0305004118000348}.
\bibitem{bkp} K. Bringmann, B. Kane and S. Pujahari,
\emph{Formulas for moments of class numbers in arithmetic
progressions}, Theorem 1.1(2).
\href{https://arxiv.org/abs/2109.07325}{arXiv:2109.07325}.
\end{thebibliography}
\end{document}
